% =====================================================================
\chapter{Use case}
\label{ch:usecase}
% =====================================================================

\section{Sơ đồ use case tổng}
Hình~\ref{fig:usecase} thể hiện 14 use case của M13: UC-01 đến UC-11 thuộc luồng
lõi được nghiệm thu; UC-12 đến UC-14 là use case hỗ trợ ở mức cần thiết. Quyền
theo H: QA Lead \emph{chạy/xem} phân tích, QC Admin \emph{cấu hình}. CVAT và
Detector là tác nhân hệ thống.

\begin{figure}[H]
\centering
\resizebox{\linewidth}{!}{%
\begin{tikzpicture}
  % ranh giới hệ thống
  \draw[lxnavy, thick, rounded corners=4pt, fill=lxgrayl!40] (-4.2,-7.6) rectangle (4.2,7.4);
  \node[font=\sffamily\small\bfseries, text=lxnavy] at (0,7.0) {LabelX — M13 Reviewer Prioritization};
  % use case cột trái
  \node[usecase] (u1)  at (-2.0, 5.8) {UC-01 Tạo snapshot từ CVAT};
  \node[usecase] (u2)  at (-2.0, 4.4) {UC-02 Chạy phân tích nghi vấn};
  \node[usecase] (u3)  at (-2.0, 3.0) {UC-03 Chấm điểm \& xếp hạng frame};
  \node[usecase] (u4)  at (-2.0, 1.6) {UC-04 Làm việc theo hàng đợi ưu tiên};
  \node[usecase] (u5)  at (-2.0, 0.2) {UC-05 Xem bằng chứng \& quyết định};
  \node[usecase] (u6)  at (-2.0,-1.2) {UC-06 Chuyển cấp \& phân xử};
  \node[usecase] (u7)  at (-2.0,-2.6) {UC-07 Rework \& kiểm lại sau sửa};
  % use case cột phải
  \node[usecase] (u8)  at ( 2.0, 5.8) {UC-08 Lập \& khoá reference};
  \node[usecase] (u9)  at ( 2.0, 4.4) {UC-09 Đo Recall@20\%};
  \node[usecase] (u10) at ( 2.0, 3.0) {UC-10 Đo effort baseline/assisted};
  \node[usecase] (u11) at ( 2.0, 1.6) {UC-11 Xem \& xuất báo cáo hiệu quả};
  \node[usecase, fill=lxamberl, draw=lxamber] (u12) at ( 2.0,-1.2) {UC-12 Tra cứu guideline};
  \node[usecase, fill=lxamberl, draw=lxamber] (u13) at ( 2.0,-2.6) {UC-13 Quản lý quyền \& audit};
  \node[font=\sffamily\scriptsize, text=lxamber] at (2.0,-3.6) {(hỗ trợ)};
  \node[usecase, fill=lxamberl, draw=lxamber] (u14) at (0,-5.4) {UC-14 Kiểm Quality Gate tối thiểu};
  % actor
  \pic at (-7.0, 1.4) {actor=Reviewer};
  \pic at (-7.0,-3.4) {actor=Annotator};
  \pic at (-7.0, 5.2) {actor={QC Admin}};
  \pic at ( 7.0, 4.6) {actor={QA Lead}};
  \pic at ( 7.0, 0.6) {actor={Product /\\Data Owner}};
  \node[ext, text width=16mm] (cvat) at (-7.0,-6.2) {«system»\\CVAT};
  \node[ext, text width=16mm] (det) at ( 7.0,-4.0) {«system»\\Detector};
  % liên kết actor
  \coordinate (rv) at (-6.6,1.8); \coordinate (an) at (-6.6,-3.0);
  \coordinate (qa) at (-6.6,5.6); \coordinate (ql) at (6.6,5.0);
  \coordinate (po) at (6.6,1.0);
  \foreach \uc in {u4,u5,u6,u7} \draw[lxnavy] (rv) -- (\uc.west);
  \draw[lxnavy] (rv) -- (u12.west);
  \draw[lxnavy] (an) -- (u7.west);
  \foreach \uc in {u1,u2,u3} \draw[lxnavy] (qa) -- (\uc.west);
  \draw[lxnavy] (qa) -- (u13.north west);
  \foreach \uc in {u6,u8,u9,u11,u14} \draw[lxnavy] (ql) -- (\uc.east);
  \draw[lxnavy] (ql) -- (u2.north east);
  \foreach \uc in {u10,u11} \draw[lxnavy] (po) -- (\uc.east);
  \draw[lxnavy] (cvat.north) -- (u1.south west);
  \draw[lxnavy] (cvat.east) -- (u7.south);
  \draw[lxnavy] (det.north) -- (u2.east);
  % include / extend
  \draw[arrd] (u2) -- node[lbl, right]{«include»} (u1);
  \draw[arrd] (u3) -- node[lbl, right]{«include»} (u2);
  \draw[arrd] (u5) -- node[lbl, right]{«include»} (u4);
  \draw[arrd] (u6) -- node[lbl, right]{«extend»} (u5);
  \draw[arrd] (u7) -- node[lbl, right]{«extend»} (u5);
  \draw[arrd] (u12) -- node[lbl, above]{«extend»} (u5);
  \draw[arrd] (u9) -- node[lbl, right]{«include»} (u8);
  \draw[arrd] (u14) -- node[lbl, left]{«include»} (u7);
\end{tikzpicture}}
\caption{Sơ đồ use case tổng của M13}
\label{fig:usecase}
\end{figure}

\section{Danh sách use case}
\begin{xltabular}{\linewidth}{@{}L{1.4cm}L{4.6cm}L{3.2cm}L{1.2cm}Y@{}}
\caption{Danh sách use case}\label{tab:uclist}\\
\toprule
\textbf{Mã} & \textbf{Tên} & \textbf{Actor chính} & \textbf{Ưu tiên} & \textbf{Mục tiêu}\\
\midrule
\endfirsthead
\toprule \textbf{Mã} & \textbf{Tên} & \textbf{Actor chính} & \textbf{Ưu tiên} & \textbf{Mục tiêu}\\ \midrule
\endhead
\bottomrule
\endfoot
UC-01 & Tạo snapshot từ CVAT & QA Lead (chạy), QC Admin (cấu hình) & \must & R-01\\
UC-02 & Chạy phân tích nghi vấn & QA Lead (chạy), QC Admin (cấu hình) & \must & R-01, R-02\\
UC-03 & Chấm điểm và xếp hạng frame & Hệ thống & \must & KPI-1\\
UC-04 & Làm việc theo hàng đợi ưu tiên & Reviewer & \must & KPI-2\\
UC-05 & Xem bằng chứng và quyết định & Reviewer & \must & R-02, KPI-2\\
UC-06 & Chuyển cấp và phân xử & QA Lead & \must & R-02\\
UC-07 & Rework và kiểm lại sau sửa & Reviewer, Annotator & \must & R-03\\
UC-08 & Lập và khoá reference & QA Lead & \must & R-04\\
UC-09 & Đo Recall@20\% & QA Lead & \must & KPI-1, R-04\\
UC-10 & Đo effort baseline/assisted & Product Owner & \must & KPI-2, R-07\\
UC-11 & Xem và xuất báo cáo hiệu quả & Product Owner, QA Lead & \must & KPI-1, KPI-2\\
UC-12 & Tra cứu guideline & Reviewer & \should & R-02\\
UC-13 & Quản lý quyền và audit & QC Admin & \must & R-05\\
UC-14 & Kiểm Quality Gate tối thiểu & QA Lead & \should & R-03, R-06\\
\end{xltabular}

\section{Đặc tả use case}
\label{sec:ucspec}

% ---------------------------------------------------------------------
\subsection*{UC-01 · Tạo snapshot từ CVAT}
\begin{xltabular}{\linewidth}{@{}>{\bfseries}L{3.2cm}Y@{}}
\toprule
Actor & QA Lead (khởi tạo); QC Admin (cấu hình phạm vi, kết nối); CVAT (hệ thống nguồn).\\
Mô tả & Đọc annotation, metadata và ảnh của phạm vi đã chọn từ CVAT, chuẩn hoá, tính hash và khoá thành snapshot bất biến.\\
Tiền điều kiện & Token CVAT được cấu hình ở backend; actor có quyền chạy phân tích trên dataset (QA Lead) theo scope; phạm vi (project/task/job) đã chọn.\\
Luồng chính &
\begin{enumerate}[nosep]
  \item Actor chọn dataset và phạm vi, bấm \emph{Tạo snapshot}.
  \item Hệ thống đọc danh sách job, frame, annotation và ảnh qua CVAT API (chỉ đọc).
  \item Hệ thống chuẩn hoá JSON annotation theo job, tính hash từng job và hash tổng.
  \item Hệ thống đọc lại dấu thời gian/hash cập nhật của từng job để kiểm drift.
  \item Không có drift: hệ thống lưu ảnh, checksum, assignee từng job, taxonomy và guideline version; khoá snapshot.
  \item Hệ thống ghi audit và hiển thị snapshot với revision hash.
\end{enumerate}\\
Luồng thay thế &
\textbf{4a.} Có drift: hệ thống huỷ khoá, báo các job đã đổi, cho phép thử lại; không tạo snapshot một phần.\newline
\textbf{2a.} Annotation không phải Bounding box: bỏ qua, ghi số lượng vào báo cáo snapshot (ngoài phạm vi).\\
Ngoại lệ & CVAT không phản hồi hoặc từ chối quyền: snapshot ở trạng thái Failed, ghi lỗi, không ảnh hưởng snapshot cũ.\newline
Actor không có quyền trong scope: backend từ chối (403), ghi audit.\\
Hậu điều kiện & Snapshot bất biến, có hash; có thể chạy UC-02.\\
Yêu cầu & \req{FR-SNP-01}…\req{FR-SNP-08}\\
\bottomrule
\end{xltabular}

% ---------------------------------------------------------------------
\subsection*{UC-02 · Chạy phân tích nghi vấn}
\begin{xltabular}{\linewidth}{@{}>{\bfseries}L{3.2cm}Y@{}}
\toprule
Actor & QA Lead (chạy); QC Admin (cấu hình engine, ngưỡng); Detector (hệ thống).\\
Mô tả & Chạy các engine MVP trên một snapshot, sinh candidate kèm evidence, cập nhật coverage ledger.\\
Tiền điều kiện & Snapshot đã khoá; cấu hình phân tích (engine, ngưỡng, seed, model artifact) có version.\\
Luồng chính &
\begin{enumerate}[nosep]
  \item Actor chọn snapshot và cấu hình, xem ước lượng thời gian và phạm vi áp dụng từng engine, bấm \emph{Chạy}.
  \item Hệ thống tạo QC Run (Queued), chia shard theo frame, đưa vào hàng đợi CPU/GPU.
  \item Worker CPU chạy Schema/Taxonomy, Geometry, Duplicate/Overlap.
  \item Worker GPU chạy Detector theo lô; worker CPU chạy matching với annotation.
  \item Mỗi shard ghi candidate, evidence và cập nhật ledger (eligible/completed units) trong cùng transaction.
  \item Khi mọi shard xong, hệ thống gộp candidate thành issue (dedup) và kích hoạt UC-03.
\end{enumerate}\\
Luồng thay thế &
\textbf{4a.} Detector lỗi trên một số frame: retry tối đa theo cấu hình; còn lỗi thì frame đó Failed cho engine này, run ở trạng thái Partial.\newline
\textbf{1a.} Không có Detector khả dụng: engine Mô hình độc lập Not checked; run vẫn chạy các engine khác.\\
Ngoại lệ & Actor huỷ run: dừng nhận shard mới, giữ kết quả đã xong, run Cancelled.\newline
Actor không có quyền chạy hoặc sửa cấu hình: backend từ chối (403).\newline
Metric engine: chỉ chạy khi có GT đã khoá cho phạm vi (B-05, B-20); thiếu GT thì Not checked, không hiển thị accuracy.\\
Hậu điều kiện & Run Completed/Partial/Failed/Cancelled; ledger phản ánh đúng phần chưa kiểm.\\
Yêu cầu & \req{FR-ENG-01}…\req{FR-ENG-10}, \req{FR-AGG-01}…\req{FR-AGG-06}\\
\bottomrule
\end{xltabular}

% ---------------------------------------------------------------------
\subsection*{UC-03 · Chấm điểm và xếp hạng frame}
\begin{xltabular}{\linewidth}{@{}>{\bfseries}L{3.2cm}Y@{}}
\toprule
Actor & Hệ thống (tự động sau UC-02); QC Admin (chọn version công thức).\\
Mô tả & Tính điểm rủi ro \(s(f)\) cho mọi frame trong snapshot, xếp hạng, lưu đóng góp từng candidate để giải thích.\\
Tiền điều kiện & Run ở trạng thái Completed hoặc Partial; công thức điểm có version đã kích hoạt.\\
Luồng chính &
\begin{enumerate}[nosep]
  \item Hệ thống tính xác suất hiệu chỉnh \(q_i\) cho từng issue theo nhóm lỗi (mục~\ref{sec:fr-rnk}).
  \item Hệ thống tính \(s(f)\) theo công thức (mục~\ref{sec:fr-rnk}), kể cả frame không có candidate (điểm nền).
  \item Hệ thống sắp xếp giảm dần, phá hoà bằng khoá cố định, lưu rank và version.
  \item Hệ thống chọn lát kiểm tra ngẫu nhiên độc lập với ranking (seed cố định).
  \item Hệ thống đánh dấu frame chưa được engine bắt buộc kiểm (Not checked/Failed) để hiển thị riêng.
\end{enumerate}\\
Luồng thay thế & \textbf{2a.} Run Partial: frame thiếu kết quả engine vẫn có điểm nhưng mang cờ ``thiếu bằng chứng''; không được coi là điểm thấp an toàn.\\
Hậu điều kiện & Mỗi frame có đúng một (score, rank, score\_version) cho run.\\
Yêu cầu & \req{FR-RNK-01}…\req{FR-RNK-09}\\
\bottomrule
\end{xltabular}

% ---------------------------------------------------------------------
\subsection*{UC-04 · Làm việc theo hàng đợi ưu tiên}
\begin{xltabular}{\linewidth}{@{}>{\bfseries}L{3.2cm}Y@{}}
\toprule
Actor & Reviewer.\\
Mô tả & Reviewer nhận frame/issue theo thứ tự ưu tiên; hệ thống cấp lease để tránh trùng việc.\\
Tiền điều kiện & Đã có ranking; reviewer có quyền Review trong scope; frame không do reviewer gán nhãn.\\
Luồng chính &
\begin{enumerate}[nosep]
  \item Reviewer mở Review Queues, chọn hàng đợi \emph{Theo rủi ro} hoặc \emph{Kiểm tra ngẫu nhiên}.
  \item Hệ thống hiển thị danh sách frame theo rank với số issue, nhóm lỗi, lý do xếp hạng.
  \item Reviewer bấm \emph{Bắt đầu review}; hệ thống cấp lease frame kế tiếp chưa có người giữ.
  \item Hệ thống mở Review Workspace (UC-05) và bắt đầu ghi effort.
\end{enumerate}\\
Luồng thay thế & \textbf{3a.} Frame kế tiếp do chính reviewer gán nhãn: hệ thống bỏ qua và cấp frame sau (self-review).\newline
\textbf{3b.} Lease hết hạn (thời hạn \TBD[-09], gia hạn khi có thao tác): frame trở về hàng đợi.\newline
\textbf{3c.} Hai reviewer cùng bấm: lease được cấp nguyên tử (một transaction, khoá dòng); người thứ hai nhận frame kế tiếp.\newline
\textbf{2a.} Hàng đợi rỗng hoặc mọi frame còn lại là self-review: hiển thị ``không còn việc hợp lệ''.\\
Hậu điều kiện & Frame ở trạng thái Đang review, có lease.\\
Yêu cầu & \req{FR-REV-01}…\req{FR-REV-04}, \req{FR-SEC-03}\\
\bottomrule
\end{xltabular}

% ---------------------------------------------------------------------
\subsection*{UC-05 · Xem bằng chứng và quyết định}
\begin{xltabular}{\linewidth}{@{}>{\bfseries}L{3.2cm}Y@{}}
\toprule
Actor & Reviewer.\\
Mô tả & Reviewer xem frame với annotation hiện tại và đề xuất của Detector, đọc evidence của từng issue, ra quyết định có lý do; có thể tạo issue mới cho lỗi không có candidate.\\
Tiền điều kiện & Reviewer giữ lease frame.\\
Luồng chính &
\begin{enumerate}[nosep]
  \item Hệ thống hiển thị ảnh, annotation (nét liền), dự đoán Detector (nét đứt), danh sách issue của frame.
  \item Reviewer chọn issue; hệ thống hiển thị evidence (lớp, confidence, IoU, rule vi phạm), guideline liên quan và case tương tự.
  \item Reviewer chọn \emph{Xác nhận lỗi}, \emph{Bác bỏ}, \emph{Chưa chắc chắn} hoặc \emph{Chuyển cấp trên}; với xác nhận chọn nhóm lỗi và mức độ.
  \item Reviewer có thể bấm \emph{Yêu cầu sửa} cho lỗi đã xác nhận (UC-07).
  \item Reviewer đánh dấu frame \emph{Đã review xong}; hệ thống lưu quyết định, actor, revision, lý do, thời gian; trả lease.
\end{enumerate}\\
Luồng thay thế & \textbf{3a.} Lỗi không có candidate: reviewer bấm vào vùng ảnh/object để tạo issue thủ công (nguồn = reviewer).\newline
\textbf{3b.} Chuyển cấp trên: issue sang UC-06.\\
Ngoại lệ & Snapshot bị thay bởi revision mới trong lúc review: quyết định vẫn gắn revision cũ; hệ thống cảnh báo.\newline
Lease đã hết hạn hoặc đã cấp cho người khác: backend từ chối lưu, yêu cầu nhận lại frame.\newline
Reviewer là assignee của job tại snapshot (mở trực tiếp qua URL): backend từ chối lưu quyết định (self-review, B-12).\\
Ghi chú KPI & Chỉ lỗi thuộc E1–E3 được đối chiếu với reference để tính KPI. Cảnh báo cấu trúc, box thừa (BR-06), lỗi độ khít biên được ghi nhận nhưng không tính KPI.\\
Hậu điều kiện & Mọi issue của frame có quyết định hoặc đang chờ phân xử; effort được ghi.\\
Yêu cầu & \req{FR-REV-05}…\req{FR-REV-12}, \req{FR-EVL-06}\\
\bottomrule
\end{xltabular}

% ---------------------------------------------------------------------
\subsection*{UC-06 · Chuyển cấp và phân xử}
\begin{xltabular}{\linewidth}{@{}>{\bfseries}L{3.2cm}Y@{}}
\toprule
Actor & QA Lead (adjudicator); Reviewer (người chuyển).\\
Mô tả & Phân xử issue chưa đủ căn cứ hoặc reviewer bất đồng; ghi Guideline Gap khi guideline thiếu.\\
Tiền điều kiện & Issue ở trạng thái Chuyển cấp trên hoặc Chưa chắc chắn.\\
Luồng chính &
\begin{enumerate}[nosep]
  \item QA Lead mở Escalations, chọn case.
  \item Hệ thống hiển thị quyết định từng reviewer, evidence, rule áp dụng.
  \item QA Lead chọn \emph{Xác nhận lỗi}, \emph{Bác bỏ} hoặc \emph{Guideline còn thiếu}, chọn nhãn đúng và rule, nhập lý do.
  \item QA Lead chọn \emph{Tạo yêu cầu sửa} hoặc \emph{Chỉ ghi nhận}.
  \item Hệ thống lưu phân xử, lưu thành Decision Case (có version) để tra cứu ở UC-12.
\end{enumerate}\\
Luồng thay thế & \textbf{3a.} Guideline còn thiếu: tạo Guideline Gap gắn rule; issue giữ trạng thái chờ đến khi guideline có version mới.\\
Hậu điều kiện & Issue có kết luận cuối hoặc gắn Guideline Gap.\\
Yêu cầu & \req{FR-ESC-01}…\req{FR-ESC-05}\\
\bottomrule
\end{xltabular}

% ---------------------------------------------------------------------
\subsection*{UC-07 · Rework và kiểm lại sau sửa}
\begin{xltabular}{\linewidth}{@{}>{\bfseries}L{3.2cm}Y@{}}
\toprule
Actor & Reviewer (yêu cầu, xác minh); Annotator (sửa); CVAT.\\
Mô tả & Lỗi đã xác nhận được sửa trên CVAT; hệ thống tạo snapshot mới, chạy lại engine bị ảnh hưởng, reviewer xác minh; cuối cùng chạy QC run cuối trên revision đã sửa (B-03).\\
Tiền điều kiện & Issue ở trạng thái Đã xác nhận.\\
Luồng chính &
\begin{enumerate}[nosep]
  \item Reviewer tạo yêu cầu sửa: việc cần làm, mức độ, hạn, annotator.
  \item Annotator mở deep link sang đúng job/frame trên CVAT, sửa, bấm \emph{Đã sửa} trên LabelX.
  \item Hệ thống tạo snapshot mới cho các job bị ảnh hưởng, chạy lại engine liên quan trên frame đã sửa (re-check).
  \item Reviewer xem kết quả re-check và frame mới, chọn \emph{Đạt} hoặc \emph{Chưa đạt}.
  \item Khi mọi yêu cầu sửa của phạm vi đã đạt, hệ thống tạo snapshot toàn phạm vi trên revision mới và QC run cuối (is\_final), liên kết các verification trước đó.
  \item Chỉ khi run cuối Completed và ledger đủ coverage engine bắt buộc, báo cáo và gate (UC-14) mới trỏ run này.
\end{enumerate}\\
Luồng thay thế & \textbf{4a.} Chưa đạt: yêu cầu mở lại (Reopened), quay về bước 2.\newline
\textbf{3a.} Re-check phát hiện lỗi mới do sửa: tạo issue mới gắn revision mới.\newline
\textbf{6a.} Run cuối Partial/Failed/Cancelled hoặc phát sinh issue mới: gate chưa đạt; issue mới vào hàng đợi; không coi phạm vi đã hoàn tất.\newline
\textbf{2a.} Người verify là annotator đã sửa: backend từ chối.\\
Hậu điều kiện & Issue Resolved chỉ khi verify trên revision đã sửa; snapshot gốc giữ nguyên.\\
Yêu cầu & \req{FR-RWK-01}…\req{FR-RWK-07}\\
\bottomrule
\end{xltabular}

% ---------------------------------------------------------------------
\subsection*{UC-08 · Lập và khoá reference}
\begin{xltabular}{\linewidth}{@{}>{\bfseries}L{3.2cm}Y@{}}
\toprule
Actor & QA Lead; hai người xác minh.\\
Mô tả & Lập GT độc lập cho tập đánh giá, suy ra tập lỗi, duyệt và khoá version (mục~\ref{sec:reference}).\\
Tiền điều kiện & Snapshot tập đánh giá đã khoá; người xác minh không có quyền xem ranking, risk score và candidate của tập này (BR-10).\\
Luồng chính &
\begin{enumerate}[nosep]
  \item QA Lead tạo phiên reference, chọn snapshot, giao hai người xác minh.
  \item Mỗi người gán GT độc lập (trong CVAT, task riêng) và nhập vào LabelX.
  \item Hệ thống matching hai bản GT, đưa phần không khớp cho QA Lead phân xử.
  \item Hệ thống suy ra \(\mathcal{E}\) theo BR-01…BR-07, hiển thị để duyệt.
  \item QA Lead duyệt; nếu sửa mapping thì hệ thống suy lại (BR-14); QA Lead khoá version.
  \item Hệ thống báo cáo tỉ lệ khớp hai bản GT trước phân xử (BR-11) và kiểm tập đánh giá không chứa case đã dùng hiệu chỉnh (BR-16).
\end{enumerate}\\
Ngoại lệ & Người xác minh là annotator của frame: hệ thống chặn giao việc.\\
Hậu điều kiện & GT và \(\mathcal{E}\) khoá version cùng mapping, ngưỡng, thuật toán.\\
Yêu cầu & \req{FR-EVL-01}…\req{FR-EVL-05}\\
\bottomrule
\end{xltabular}

% ---------------------------------------------------------------------
\subsection*{UC-09 · Đo Recall@20\%}
\begin{xltabular}{\linewidth}{@{}>{\bfseries}L{3.2cm}Y@{}}
\toprule
Actor & QA Lead.\\
Mô tả & Tính KPI-1 trên tập held-out, so với ranking đối chứng, kèm khoảng tin cậy.\\
Tiền điều kiện & Reference khoá; ranking của held-out đã tính bằng version công thức đã khoá trên tập hiệu chỉnh.\\
Luồng chính &
\begin{enumerate}[nosep]
  \item QA Lead chọn run, reference version, \(k\) (mặc định 20\%).
  \item Hệ thống kiểm \(|\mathcal{E}| \ge E_{min}\); tính Recall@k, recall theo nhóm lỗi, theo slice, recall theo frame có lỗi.
  \item Hệ thống tính các đường đối chứng (ngẫu nhiên nhiều seed, heuristic) và khoảng tin cậy bootstrap.
  \item Hệ thống lưu evaluation run với đầy đủ provenance.
\end{enumerate}\\
Luồng thay thế & \textbf{2a.} \(|\mathcal{E}| < E_{min}\): ghi ``không đủ mẫu'', không kết luận.\\
Quy tắc & Chỉ đếm lỗi E1–E3 theo BR-01…BR-07; loại đối tượng ignore khỏi tử số và mẫu số; box thừa, cảnh báo cấu trúc và độ khít biên báo cáo riêng.\newline
Recall@20\% đo khả năng xếp hạng, không thay thế phép đo accuracy của engine. Accuracy của Detector (Metric engine) đo riêng trên GT đã khoá; Model Orchestrator tổng hợp kết quả có provenance (B-05).\\
Yêu cầu & \req{FR-EVL-07}…\req{FR-EVL-10}, \req{FR-EVL-15}\\
\bottomrule
\end{xltabular}

% ---------------------------------------------------------------------
\subsection*{UC-10 · Đo effort baseline/assisted}
\begin{xltabular}{\linewidth}{@{}>{\bfseries}L{3.2cm}Y@{}}
\toprule
Actor & Product Owner (thiết kế thí nghiệm); Reviewer (tham gia); QA Lead (giám sát).\\
Mô tả & Chạy thí nghiệm hai nhánh theo thiết kế ở mục~\ref{sec:kpi2}, ghi effort tự động, tính KPI-2 và kiểm định chất lượng tương đương.\\
Tiền điều kiện & Thiết kế thí nghiệm (phân nhóm, quy tắc dừng, \(\delta\)) đã khoá; reference của dữ liệu thí nghiệm đã khoá.\\
Luồng chính &
\begin{enumerate}[nosep]
  \item Product Owner tạo thí nghiệm, gán reviewer vào nhánh theo thiết kế chéo.
  \item Nhánh baseline: Workspace ẩn ranking và evidence, frame theo thứ tự gốc.
  \item Nhánh assisted: Workspace hiển thị ranking và evidence, dừng theo quy tắc đã khoá.
  \item Hệ thống ghi effort theo hoạt động (review, cảnh báo sai, phân xử, re-check).
  \item Kết thúc: hệ thống tính \(T_A, T_B, \rho_A, \rho_B\), KPI-2, KPI-2b và kiểm định non-inferiority.
\end{enumerate}\\
Yêu cầu & \req{FR-EVL-11}…\req{FR-EVL-16}\\
\bottomrule
\end{xltabular}

% ---------------------------------------------------------------------
\subsection*{UC-11 · Xem và xuất báo cáo hiệu quả}
\begin{xltabular}{\linewidth}{@{}>{\bfseries}L{3.2cm}Y@{}}
\toprule
Actor & Product Owner, QA Lead.\\
Mô tả & Xem báo cáo KPI-1, KPI-2, guardrail, coverage, có ghi cỡ mẫu và phương pháp; xuất PDF/CSV/JSON.\\
Luồng chính &
\begin{enumerate}[nosep]
  \item Actor chọn evaluation run/thí nghiệm.
  \item Hệ thống hiển thị: đường Recall@k, bảng theo nhóm lỗi và slice, effort hai nhánh, residual, guardrail, coverage từng engine.
  \item Actor xuất báo cáo; tệp xuất ghi version của snapshot, run, công thức, reference.
\end{enumerate}\\
Quy tắc & Chỉ số chưa đủ dữ liệu hiển thị Not checked/không đủ mẫu, không hiển thị 0\% hay đạt.\\
Yêu cầu & \req{FR-RPT-01}…\req{FR-RPT-06}\\
\bottomrule
\end{xltabular}

% ---------------------------------------------------------------------
\subsection*{UC-12 · Tra cứu guideline (hỗ trợ)}
\begin{xltabular}{\linewidth}{@{}>{\bfseries}L{3.2cm}Y@{}}
\toprule
Actor & Reviewer.\\
Mô tả & Từ issue, xem rule guideline liên quan theo rule ID + version + section và các Decision Case đã duyệt (B-13); không dùng RAG.\\
Luồng chính & Hệ thống tra mapping (nhóm lỗi, lớp) $\rightarrow$ rule ID; hiển thị trích đoạn, version, section; liệt kê case đã phân xử cùng rule.\\
Yêu cầu & \req{FR-GDL-01}…\req{FR-GDL-04}\\
\bottomrule
\end{xltabular}

% ---------------------------------------------------------------------
\subsection*{UC-13 · Quản lý quyền và audit (hỗ trợ)}
\begin{xltabular}{\linewidth}{@{}>{\bfseries}L{3.2cm}Y@{}}
\toprule
Actor & QC Admin; Super Admin (ghi đè có lý do).\\
Mô tả & Gán vai trò theo scope dataset; xem nhật ký kiểm toán append-only.\\
Quy tắc & Backend luôn kiểm quyền; ẩn nút chỉ để thuận tiện. Người duyệt khác người yêu cầu. Ghi đè của Super Admin phải có lý do và gắn nhãn trong audit.\\
Yêu cầu & \req{FR-SEC-01}…\req{FR-SEC-07}\\
\bottomrule
\end{xltabular}

% ---------------------------------------------------------------------
\subsection*{UC-14 · Kiểm Quality Gate tối thiểu (hỗ trợ)}
\begin{xltabular}{\linewidth}{@{}>{\bfseries}L{3.2cm}Y@{}}
\toprule
Actor & QA Lead (xem, đề nghị waiver); người có quyền duyệt ngoại lệ (khác người đề nghị).\\
Mô tả & Tính từng điều kiện gate trên QC run cuối; hiển thị nguồn, cỡ mẫu, kết quả; xử lý waiver theo policy. Không bao gồm phát hành dataset đầy đủ.\\
Tiền điều kiện & Có QC run cuối (UC-07) hoặc run gốc nếu chưa có rework (ghi rõ).\\
Luồng chính &
\begin{enumerate}[nosep]
  \item Hệ thống tính từng điều kiện theo FR-GTE-01, mỗi điều kiện có nguồn, mẫu số, cỡ mẫu.
  \item Điều kiện thiếu dữ liệu hoặc engine bắt buộc Not checked/Partial/Failed: \emph{chưa đạt}.
  \item QA Lead xem kết quả; nếu có điều kiện chưa đạt và policy cho phép, tạo đề nghị waiver với lý do, evidence, hạn.
  \item Người duyệt khác người đề nghị duyệt/từ chối; waiver chỉ hiệu lực sau duyệt.
\end{enumerate}\\
Ngoại lệ & Người duyệt trùng người đề nghị: backend từ chối. Waiver hết hạn: điều kiện trở lại chưa đạt.\\
Yêu cầu & \req{FR-GTE-01}…\req{FR-GTE-04}\\
\bottomrule
\end{xltabular}
